\section{A smoothness obstruction for finite flat torsors}

The infinitesimal Picard kernel in Theorem~\ref{thm:surface} does not by itself give a nontrivial point of that kernel. The following obstruction explains why one natural finite flat torsor construction cannot bridge the difference while preserving smoothness.

\begin{proposition}\label{prop:tangent-ample}
Let $Z$ be a smooth projective surface over an algebraically closed field and let $D\subset Z$ be a smooth ample integral curve. There is no line subbundle
\[
 \OO_Z(-D)\longrightarrow T_Z
\]
whose restriction to $D$ is tangent to $D$. Here a subbundle has locally free quotient. No hypothesis of $p$-closedness is imposed.
\end{proposition}

\begin{proof}
Write $\mathcal F$ for the proposed subbundle. Its restriction maps into $T_D$ and is nonzero on every fibre. Thus
\[
 \OO_D(-D)\simeq T_D,\qquad D^2=2g(D)-2,\qquad K_Z\cdot D=0.
\]
The quotient of $T_Z$ by $\mathcal F$ has first Chern class $D-K_Z$. Multiplicativity of Chern classes gives
\begin{equation}\label{eq:tangent-chern}
 c_2(Z)=(-D)\cdot(D-K_Z)=-D^2<0.
\end{equation}

In characteristic zero one can already contradict ampleness on the curve. A rank-one distribution is involutive, and its bracket induces a regular connection on $N_{D/Z}$: for a tangent field $v$ and a normal class represented by $w$, put $\nabla_v\overline w=\overline{[v,w]}$. Tangency makes this independent of the lifts, and the bracket identities give the connection rules. A line bundle with a regular connection on a smooth projective curve has degree zero in characteristic zero. Indeed, in a meromorphic frame the connection is a meromorphic differential whose residue at each point is the order of that frame there; the residue theorem says that their sum, the degree, is zero. This contradicts $\deg N_{D/Z}=D^2>0$.

Suppose now that the characteristic is positive. If a positive multiple of $K_Z$ has a nonzero section, its effective divisor has zero intersection with the ample divisor $D$. It is therefore zero. The canonical bundle is torsion, $Z$ is minimal, and its Kodaira dimension is zero. The classification gives $c_2(Z)\in\{0,12,24\}$, including the nonclassical cases in characteristics two and three; see \cite[\S7, author pp.~26--30]{Liedtke}. This contradicts \eqref{eq:tangent-chern}.

Consequently the Kodaira dimension is $-\infty$. The surface is birationally ruled by \cite[Theorems~4.4--4.5, author p.~21]{Liedtke}. Rational surfaces have positive second Chern number, so a minimal model is a ruled surface $Z_0\to C$ of genus $g\geq2$. Let $r$ be the number of point blowups from $Z_0$ to $Z$. Choose a section class $C_0$ with $C_0^2=-e$, a fibre class $F$, and total-transform exceptional classes $E_1,\ldots,E_r$. Pulling back the first two classes to $Z$, write
\[
 D=aC_0+bF-\sum_i m_iE_i,\qquad
 K_Z=-2C_0+(2g-2-e)F+\sum_iE_i.
\]
The intersection form has $C_0\cdot F=1$, $F^2=0$, and $E_i\cdot E_j=-\delta_{ij}$; each $E_i$ is orthogonal to the pullback of the minimal model. Moreover
\begin{equation}\label{eq:exceptional-bounds}
 a\geq1,\qquad 0<m_i<a.
\end{equation}
To check the strict upper bound even at infinitely near centres, let $F_c$ be the full fibre containing the centre that creates $E_i$. At that blowup its total transform contains the exceptional divisor with positive multiplicity. Subtracting one copy leaves a nonzero effective divisor, and subsequent pullbacks preserve this property. Thus both $E_i$ and $F_c-E_i$ are nonzero effective divisors. Their intersections with $D$ are respectively $m_i$ and $a-m_i$, and ampleness proves \eqref{eq:exceptional-bounds}.

Now $K_Z\cdot D=0$ and the intersection form give
\[
 2b=a(e+2g-2)+\sum_i m_i,\qquad
 D^2=2a^2(g-1)+\sum_i m_i(a-m_i).
\]
On the other hand $c_2(Z)=4(1-g)+r$. Equation~\eqref{eq:tangent-chern} therefore becomes
\[
 (2a^2-4)(g-1)+\sum_i\bigl(m_i(a-m_i)+1\bigr)=0.
\]
For $a\geq2$ its left side is strictly positive. For $a=1$, the strict bounds on the integers $m_i$ force $r=0$, and the equation forces $g=1$. Both alternatives contradict the preceding conclusions.
\end{proof}

Regular foliations in positive characteristic can behave differently from their characteristic-zero counterparts. Fassarella, Mendson, dos Santos and Touzet classify them on minimal rational and del Pezzo surfaces and prove an arbitrary-characteristic Camacho--Sad formula \cite[Theorems~1.1--1.2 and 2.8]{FMST}. Their results provide the surrounding foliation theory. Proposition~\ref{prop:tangent-ample} uses the particular tangent line bundle and also excludes the ruled surfaces of higher genus forced by its negative Chern number.

\begin{corollary}\label{cor:smooth-torsor}
Let $X$ be a smooth projective surface over an algebraically closed field of characteristic $p>0$, and let $Y\subset X$ be a smooth ample integral curve. Put $L=\OO_X(Y)$ and let $s$ be its canonical section. Define the finite flat group scheme $G_s$ locally by
\[
 \OO_X[u]/(u^p),\qquad
 \Delta(u)=u\otimes1+1\otimes u+s\,u\otimes u.
\]
If $s_j=g_{ij}^{-1}s_i$ in two trivializations of $L$, use $u_j=g_{ij}u_i$ to glue these algebras. There is no $G_s$-torsor $\pi:Z\to X$ with $Z$ smooth and the full inverse image $D=\pi^*Y$ smooth and integral.
\end{corollary}

\begin{proof}
The formula defines a rank-$p$ group scheme with $\operatorname{Lie}(G_s)=L^{-1}$. The homomorphism to $\mu_p$ given by $1+su$ is an isomorphism off $Y$; over $Y$ the group is the additive height-one group with the indicated line twist.

A torsor is finite flat over $X$, so $Z$ is projective and $D$ is ample. Differentiating the action gives
\[
 \pi^*\operatorname{Lie}(G_s)=\OO_Z(-D)\longrightarrow T_Z.
\]
At every geometric point the orbit map identifies the fibre group with the fibre of the torsor through that point, a closed subscheme of $Z$. Its tangent map is injective. Hence the displayed map has rank one on every fibre and has locally free quotient. Its derivations annihilate pullbacks of functions from $X$, so the image is tangent to $D$. Proposition~\ref{prop:tangent-ample} is the required contradiction.
\end{proof}

Resolving a singular torsor does not preserve these hypotheses: the pullback of an ample curve can become only nef and big, and the tangent subsheaf can acquire zeros. The corollary therefore leaves the general kernel on geometric Picard points undecided.
